\section{Strategie ewolucyjne}

Strategie ewolucyjne (ang. \textit{evolution strategies} --- ES) to gałąź algorytmów ewolucyjnych operująca na wektorach liczb rzeczywistych, a algorytmy te dobrze sprawdzają się w problemach optymalizacji ciągłej w przestrzeniach o dużej wymiarowości \cite{Hansen2016CMA}. Strategia ta operuje na populacji wektorów $\mathbf{x} \in \mathbb{R}^n$, a głównym operatorem mutacji jest \textbf{mutacja gaussowska}:
\[
x_i' = x_i + \mathcal{N}(0, \sigma_i^2)
\]
gdzie $\sigma_i$ to odchylenie standardowe, zwane \textit{wielkością kroku} (ang. \textit{step size}). W wielu wariantach strategii ewolucyjnych stosuje się \textbf{samodostosowywanie} parametrów mutacji --- wektory $\boldsymbol{\sigma}$ mogą być aktualizowane równolegle z samym rozwiązaniem, co ułatwia dostosowanie skali przeszukiwań do lokalnej geometrii przestrzeni funkcji celu \cite{Hansen2016CMA}, jednak nie jest to cecha wspólna wszystkich algorytmów z tej rodziny.

Strategie ewolucyjne określa się za pomocą notacji $(\mu + \lambda)$ lub $(\mu, \lambda)$, które opisują ich dwa różne warianty, gdzie $\mu$ oznacza liczbę rodziców, a $\lambda$ liczbę potomków. W wariancie $(\mu + \lambda)$ selekcja obejmuje zarówno rodziców, jak i potomstwo (podejście elitarne), natomiast w wariancie $(\mu, \lambda)$ populacja rodziców jest w całości zastępowana przez potomstwo, co sprzyja eksploracji \cite{Hansen2016CMA}.

Strategie ewolucyjne zastosowano m.in. w pracy García-Sáncheza i in.~\cite{sanchez2020} do optymalizacji wag funkcji oceny agenta grającego w Hearthstone.

